\section{Evaluation Strategy}
\label{sec:evaluation-strategy}

The evaluation of this thesis addresses three distinct questions that together cover the four research questions of \cref{sec:final-problem-statement}: whether the system functions correctly as a platform for reading experiments, whether the architecture satisfies the modularity and extensibility goals that motivate the design, and whether its runtime behaviour meets the live-operation constraints that adaptive reading imposes.
Tab.~\ref{tab:rq-evaluation} maps each research question to the method by which it is evaluated, and the remainder of this section describes those methods.

\begin{table}[htbp]
\centering
\caption{Mapping of each research question to its evaluation method. The reading-behaviour effect study that a full evaluation of intervention efficacy would require is outside the scope of this thesis.}
\label{tab:rq-evaluation}
\begin{tabular}{p{0.30\linewidth} p{0.60\linewidth}}
\toprule
Research question & Evaluation method \\
\midrule
RQ1 --- Modularity & Dependency-graph inspection (no imports across module boundaries), module-installer review, and demonstration that a reference provider connects without modifying the core. \\
RQ2 --- Sensing pipeline & Measurement of end-to-end gaze-to-event and gaze-to-intervention latency under live operation, and confirmation that the analysis stage is inspectable, loggable, and replayable. \\
RQ3 --- Context-preserving intervention & Demonstration that interventions are committed at controlled boundaries, together with inspection of the logged context-preservation events. \\
RQ4 --- Researcher control & Demonstration of the researcher's session controls and verification of the auditable record through session replay and data export. \\
\bottomrule
\end{tabular}
\end{table}

For functional correctness, we evaluate the system against the functional requirements defined in \cref{ch:requirements}, verifying that each requirement is met through a combination of automated tests, manual end-to-end walkthroughs of the researcher and participant workflows, and inspection of exported session data.
For architectural quality, we evaluate the degree to which the modular boundaries between sensing, decision-making, intervention execution, and user interface hold in the implemented system, and whether the external module provider integration contract is sufficient for a third-party provider to connect without modifying the core application.
Concretely, the architectural-quality assessment summarised in the RQ1 row of Tab.~\ref{tab:rq-evaluation} is carried out through three means:
\begin{enumerate}
  \item Inspection of the dependency graph to confirm that no module imports across the boundaries defined in the design.
  \item Demonstration that the external integration contract is sufficient for the reference decision-maker sub-project to connect to the core runtime without modifying it.
  \item Review of the module installer pattern to confirm that new implementations can be substituted at the composition root without touching application logic.
\end{enumerate}

The remaining methods in Tab.~\ref{tab:rq-evaluation} concern the runtime behaviour of the platform rather than its structure: the gaze-to-event and gaze-to-intervention latency measurement for RQ2, and the boundary-commit and session-replay demonstrations for RQ3 and RQ4, are carried out as part of the evaluation reported in \cref{ch:evaluation}.

We acknowledge one structural threat to validity: as a two-person implementation team, we cannot assess our own architectural decisions with full objectivity.
We partially mitigate this, by enabling external scrutiny: we publish the codebase and integration contracts openly under the MIT License, making them available for independent inspection, and we treat the reference decision-maker and reference eye-movement analyser as independently exercised validators of the external integration contract rather than as internal tests.

We do not conduct a user study of reading behaviour within the scope of this thesis.
The platform is designed to enable such studies, but evaluating the effectiveness of specific typographic interventions on reading comprehension or fatigue is a separate research question that requires a dedicated experiment with recruited participants and an approved study protocol.
That question is explicitly outside the scope of this work, as discussed in \cref{ch:introduction}.
